\section{Supplementary material}

\subsection{Reproducing the study}\label{sec:repro}
Everything below is in the project repository. Run in this order.

\begin{enumerate}
\item \textbf{Environment.} Python~3.11 with ECG-Image-Kit's core requirements (handwriting extras omitted); Python~3.13+ with NumPy, SciPy, pandas, Matplotlib, OpenCV, \texttt{wfdb}, NeuroKit2 for analysis; Google Colab (T4) for digitisation. \emph{Record the commit of every tool at the start of each run} (not done in this study).
\item \textbf{Sweep 1.} \texttt{colab/MediSense\_Cardio\_FullSweep\_v1.ipynb} (built by \texttt{scripts/notebook\_builders/build\_fullsweep\_notebook.py}). Set the runtime to T4 GPU, confirm the log prints \texttt{Using device: cuda}, and run all; it resumes from \texttt{results.csv} on Drive.
\item \textbf{Sweep 1 analysis.} \texttt{colab/MediSense\_Cardio\_Analysis\_Part1.ipynb} (cleaning), \texttt{Part1b} (failing records), \texttt{Part2} (heart-rate outcomes), \texttt{Part3\_figures}.
\item \textbf{Sweep 2.} \texttt{colab/MediSense\_Cardio\_FullSweep\_v2\_Intervals.ipynb}, writing \texttt{results\_v2.csv}.
\item \textbf{Sweep 2 analysis (local).} \texttt{scripts/analysis/analyze\_v2\_intervals.py} (cleaning, intervals, QTc outcome), \texttt{hr\_outcomes\_v2.py}, \texttt{check\_gt\_sampling\_rate.py} and \texttt{redo\_intervals\_gt500.py} (500\,Hz ground truth), \texttt{journal\_extra\_analysis.py}, \texttt{journal\_tables.py}, \texttt{journal\_latex\_tables.py}.
\item \textbf{Figures.} \texttt{journal\_fig\_degradation\_grid.py}, \texttt{journal\_figs\_results.py}; Sweep~1 figures from \texttt{Part3\_figures}.
\end{enumerate}

\subsection{Results file data dictionary}
\begin{longtable}{@{}>{\ttfamily}p{3.6cm}p{11.2cm}@{}}
\caption{Columns of the Sweep~2 results file (\texttt{results\_v2.csv}, 1{,}550 rows).}\label{tab:dict}\\
\toprule
Column & Meaning \\
\midrule
\endfirsthead
variant\_id & Unique test identifier: record, axis and level (for example \texttt{15049\_lr-0\_\_blur5}). \\
record & PTB-XL record name. \\
axis & \texttt{clean}, \texttt{blur}, \texttt{rotate}, \texttt{jpeg}, \texttt{resolution}, \texttt{brightness} or \texttt{gridcontrast}. \\
level & 0 (clean) or 1--5. \\
sex & \texttt{male} or \texttt{female}, from PTB-XL metadata. \\
gt\_hr & Ground-truth lead~II heart rate (bpm), from the 100\,Hz record. \\
gt\_pr\_ms, gt\_qrs\_ms, gt\_qt\_ms, gt\_qtc\_ms & Ground-truth intervals (ms) from the 100\,Hz record (the analysis recomputes them at 500\,Hz). \\
gt\_qtc\_abnormal & Ground-truth QTc above the sex-specific threshold. \\
dig\_hr, dig\_pr\_ms, dig\_qrs\_ms, dig\_qt\_ms, dig\_qtc\_ms & Values from the digitised signal; empty if none could be extracted. \\
dig\_qtc\_abnormal & Digitised QTc above the sex-specific threshold. \\
diff\_bpm, diff\_pr\_ms, diff\_qrs\_ms, diff\_qt\_ms, diff\_qtc\_ms & Absolute difference between digitised and ground-truth values. \\
snr\_db & \textbf{Invalid; do not use} (Section~\ref{sec:metrics}). \\
\bottomrule
\end{longtable}

\subsection{Digitiser configuration}
Key settings passed to Open-ECG-Digitizer in both sweeps: segmentation U-Net weights \texttt{unet\_weights\_07072025.pt} (4 output channels, 8 stages, depth 2); lead-name U-Net \texttt{lead\_name\_unet\_weights\_07072025.pt}; layout file \texttt{lead\_layouts\_george-moody-2024.yml}; \texttt{target\_num\_samples} 10{,}000; \texttt{resample\_size} 3{,}000; dewarping off; perspective detector 250 angles; pixel-size finder 30--70 grid lines. The full generated configuration is in the sweep notebooks.

\subsection{Corrections made while preparing this manuscript}
For transparency, the following were found and corrected during preparation, and are recorded so that readers of earlier drafts know what changed.
\begin{itemize}
\item The planned SNR tier was found to be computed on unaligned signals and was removed.
\item A first attribution of the interval error floor to delineation alone was revised after the 500\,Hz ground-truth check.
\item A first QTc false-reassurance count of 39 (100\,Hz ground truth) or 26 (500\,Hz) was found to be dominated by records already misclassified on their clean image, and was corrected to 6.
\item An earlier draft labelled a mean as a 90th percentile; several citations were re-assigned after re-checking the reference list against the sources.
\item A patient age of 300 in PTB-XL was recognised as a de-identification code.
\end{itemize}
